\documentclass{situationist-preprint}

\title{The Open Future}
\msid{TSM-12-a}
\msfamily{TSM-12}
\msdiscipline{Time, Space and Modality}
\mstype{Principal manuscript}
\mskeywords{open future; future contingents; branching time; settledness; thin red line; B-theory}

\begin{document}
\maketitle

\begin{abstract}
We argue that the future is open: future contingents are neither true nor false until they are settled. What has a truth value at a moment has it in virtue of what is the case relative to that moment: its place in the branching structure, its past and the laws. When a sentence is true on some histories through a moment and false on others, nothing relative to the moment decides between them; a \enquote{thin red line} would be an ungrounded mark privileging one history as an objective present would privilege one time. The thesis fits a tenseless B-theory: settledness is moment-relative, every moment has a settled past and an open future relative to itself, and no moment is the edge of a growing block.
\end{abstract}

\section{The problem}

Aristotle asked whether it is already true that there will be a sea battle tomorrow, and on the common reading answered that it is not yet true or false (\emph{De Interpretatione} 9). \citet{Thomason1970} kept classical logic while denying bivalence by supervaluating over the histories through a moment. B-theorists \citep{Williams1951,Sider2001} commonly hold that later times are as determinate as earlier ones; growing-block theorists \citep{Broad1923} secure openness by denying that the future exists, at the price of a present edge of reality.

We defend a third position: the future is open within a tenseless branching structure in which no moment is privileged. ACT-17 states the semantics and its consequence for fatalism, assuming that free choices are open; here we give the independent ground on which it relies.

\section{Framework}

\begin{definition}[Branching structure]\label{d:branch}
\emph{Moments} are partially ordered by temporal precedence. The order branches toward the later but not toward the earlier: the predecessors of any moment are linearly ordered, and any two moments have a common predecessor. A \emph{history} is a maximal linearly ordered set of moments; the histories through $m$ are its possible continuations, and the \emph{past} of $m$ is the set of its predecessors.
\end{definition}

The definition is that of ACT-17-a. The structure exists tenselessly, and every history through $m$ contains the whole past of $m$.

\begin{definition}[Settledness]\label{d:settled}
A sentence is \emph{settled true} (\emph{false}) \emph{at} $m$ if it is true (false) relative to $m$ and every history through $m$, and \emph{open at} $m$ otherwise. A \emph{future contingent at} $m$ is a sentence about what happens after $m$ that is open at $m$.
\end{definition}

\begin{principle}[Moment-relative grounding]\label{p:ground}
If a sentence has a truth value at a moment $m$, it has it in virtue of what is the case relative to $m$: the branching structure, the position of $m$ in it, the course of events up to and at $m$, and the laws.
\end{principle}

\begin{principle}[No thin red line]\label{p:trl}
Nothing that is the case relative to $m$ distinguishes one history through $m$ from the others as the one that will be followed.
\end{principle}

\begin{assumption}[Contingency]\label{a:cont}
There are moments $m$ and sentences about what happens after $m$ that are true relative to some histories through $m$ and false relative to others.
\end{assumption}

Principle~\ref{p:ground} is a moment-relative form of the thesis that truth supervenes on being \citep{Bigelow1988}: a value at a moment is a value from its standpoint, which only what that standpoint contains can ground. Principle~\ref{p:trl} is defended in Objection~\ref{o:trl}. Assumption~\ref{a:cont} is indeterminism; ACT-17 locates its paradigm cases in free choices.

\section{The argument}

\begin{lemma}\label{l:true}
If a sentence is true (false) at $m$, it is true (false) relative to every history through $m$.
\end{lemma}

\begin{proof}
By Principle~\ref{p:ground}, a truth at $m$ is grounded in what is the case relative to $m$, which by Principle~\ref{p:trl} is the same whichever history through $m$ is considered; so it makes $\phi$ true relative to each. Falsity is parallel.
\end{proof}

\begin{theorem}\label{t:main}
A future contingent at $m$ is neither true nor false at $m$. Since there are future contingents, the future is open.
\end{theorem}

\begin{proof}
A future contingent at $m$ is true relative to some history through $m$ and false relative to another (Definition~\ref{d:settled}). By Lemma~\ref{l:true}, it is neither true nor false at $m$. By Assumption~\ref{a:cont}, there are such sentences.
\end{proof}

\begin{proposition}\label{p:settle}
A sentence open at $m$ is settled at every later moment $m'$ by which the matter has occurred, and no sentence about the past of a moment is open at that moment.
\end{proposition}

\begin{proof}
Every history through $m'$ contains the past of $m'$, so an event in it, or at $m'$, is on all of them or on none.
\end{proof}

Relative to every moment, then, the past is fixed and the contingent future open; the asymmetry lies in the shape of the structure, not in a privileged present.

\section{Compatibility with the B-theory}\label{s:btheory}

The view respects TSM-01, TSM-15 and TSM-16: branching is a feature of the B-series, and nothing is objectively present. \enquote{Settled at $m$}, like \enquote{north of $p$}, is a relation in which every moment stands to some sentences; to say that a contingent is open \enquote{now} is to say it is open at the moment of utterance (TSM-18). This is not the growing block: existence is not relativized to moments, only settledness is, and there is no edge whose inhabitants could not know they were at it \citep{Bourne2002,BraddonMitchell2004}.

Tenseless sentences are evaluated at the actual history (ACT-17-a); but \enquote{actual}, like \enquote{now}, is indexical \citep{Lewis1970}, and from a moment's standpoint the actual history is fixed only as far as that moment's past. Nor do we claim that the future of a moment exists as a single region: relative to $m$ the past is a unique chain, while the continuations are many possibilities, not concrete realms (compare TSM-08). Openness is unsettledness of sentences at a moment, not a state of affairs neither obtaining nor failing to obtain (MET-19); each history is determinate.

\section{Objections and replies}

\begin{objection}[The thin red line]\label{o:trl}
One history through $m$ is the actual future, and \enquote{it will be that $p$} is true at $m$ if $p$ happens on it \citep{MalpassWawer2012}.
\end{objection}

\begin{reply}
A fact singling out one history through $m$ is either grounded in what happens up to $m$ and the laws, or not. If it is, the other histories are not genuine continuations and the sentence is not contingent. If not, it is a primitive mark on one branch, since everything occurring at $m$ is common to the histories through it; and \citet{BelnapGreen1994} argued that no assignment of such marks fits the interaction of tense with possibility. A primitive mark on one branch is, moreover, to the branching structure what an objective present is to the B-series, and TSM-15 and TSM-16 reject the latter.
\end{reply}

\begin{objection}[Retrospective truth]\label{o:retro}
After the battle we say that it was true yesterday that there would be a battle.
\end{objection}

\begin{reply}
At the later moment \enquote{it was the case that there will be a battle} is settled true \citep{Thomason1970}; it does not follow that \enquote{there will be a battle} was true at the earlier one \citep[compare][]{MacFarlane2003}.
\end{reply}

\begin{objection}[The determinate future]\label{o:det}
On the B-theory every later event is tenselessly located, so future-tensed sentences are already true or false \citep{Williams1951,Sider2001}.
\end{objection}

\begin{reply}
The objection assumes a linear order of moments, which the B-theory does not require. Whether time branches is a question about the laws (Assumption~\ref{a:cont}), not about tense.
\end{reply}

\section{Conclusion}

Nothing relative to a moment decides a later contingent matter, so future contingents are neither true nor false until settled. The future is open in a tenseless branching structure with neither a present edge nor a privileged branch.

\begin{thebibliography}{99}\small
\bibitem[Belnap and Green(1994)]{BelnapGreen1994} Belnap, N. and Green, M. (1994). Indeterminism and the thin red line. \emph{Philosophical Perspectives}, 8, 365--388.
\bibitem[Bigelow(1988)]{Bigelow1988} Bigelow, J. (1988). \emph{The Reality of Numbers: A Physicalist's Philosophy of Mathematics}. Clarendon Press.
\bibitem[Bourne(2002)]{Bourne2002} Bourne, C. (2002). When am I? A tense time for some tense theorists? \emph{Australasian Journal of Philosophy}, 80(3), 359--371.
\bibitem[Braddon-Mitchell(2004)]{BraddonMitchell2004} Braddon-Mitchell, D. (2004). How do we know it is now now? \emph{Analysis}, 64(3), 199--203.
\bibitem[Broad(1923)]{Broad1923} Broad, C.~D. (1923). \emph{Scientific Thought}. Kegan Paul.
\bibitem[Lewis(1970)]{Lewis1970} Lewis, D. (1970). Anselm and actuality. \emph{No\^us}, 4(2), 175--188.
\bibitem[MacFarlane(2003)]{MacFarlane2003} MacFarlane, J. (2003). Future contingents and relative truth. \emph{The Philosophical Quarterly}, 53(212), 321--336.
\bibitem[Malpass and Wawer(2012)]{MalpassWawer2012} Malpass, A. and Wawer, J. (2012). A future for the thin red line. \emph{Synthese}, 188(1), 117--142.
\bibitem[Sider(2001)]{Sider2001} Sider, T. (2001). \emph{Four-Dimensionalism: An Ontology of Persistence and Time}. Oxford University Press.
\bibitem[Thomason(1970)]{Thomason1970} Thomason, R.~H. (1970). Indeterminist time and truth-value gaps. \emph{Theoria}, 36(3), 264--281.
\bibitem[Williams(1951)]{Williams1951} Williams, D.~C. (1951). The myth of passage. \emph{The Journal of Philosophy}, 48(15), 457--472.
\end{thebibliography}

\end{document}
